\GuideSpread{The radical keeps one copy of each prime}{Support and repeated depth are different measurements}{01_radical.pdf}{
The radical of a positive integer is the product of its distinct prime divisors, with $\rad(1)=1$. In $72=2^3 3^2$, the radical is $6$. The shaded first copies mark support. The remaining copies mark repeated depth.\par
For a primitive triple $a+b=c$, let $R=\rad(abc)$. Pairwise coprimality makes the three prime sectors disjoint, so $R=R_aR_bR_c$. The conjecture compares the largest term $c$ with this support, allowing every positive $\varepsilon$.\par
It does not bound the exponents individually. The additive equation is the mechanism that must constrain their simultaneous growth.
}{Retain the distinction between activated prime support and prime power multiplicity.}

\GuideSpread{A complete example fits in the margin}{The primitive triple $1+8=9$}{02_small_triple.pdf}{
The inputs are $1$ and $8=2^3$. The output is $9=3^2$. Their supports are disjoint, and $R=2\cdot3=6$. Thus $c/R=3/2$.\par
The standard quality is $q=\log c/\log R$, here approximately $1.226294$. A quality above one means $c>R$. It does not contradict abc. The conjecture permits a constant depending on $\varepsilon$ and allows finitely many exceptional triples at each fixed margin.\par
The total repeated part is $Q=abc/R=12$. Dividing by the input scale $ab=8$ gives the balanced ratio $3/2$. That cancellation will matter throughout the guide.
}{Use this triple to check every normalization by hand.}

\GuideSpread{High quality is an exact arithmetic feature}{The worked triple $2+109\cdot3^{10}=23^5$}{03_high_quality.pdf}{
The exact output is $6{,}436{,}343$. Its radical is $R=2\cdot3\cdot109\cdot23=15{,}042$, and the quality is approximately $1.629912$. The factor $109$ is essential to the equality.\par
Repeated powers occur in two sectors. The input has ten copies of $3$, while the output has five copies of $23$. The radical counts each of these primes once.\par
This example demonstrates what a uniform theorem must accommodate. It does not set a universal quality ceiling, prove an optimal constant, or belong to the later census with $c\le10{,}000$. Its equality, coprimality, and factorization were checked exactly.
}{An example can test a proposed bound without establishing a universal bound.}

\GuideSpread{A bounded census has a precise domain}{All quality above one triples with $a\le b$ and $c\le10{,}000$}{04_bounded_census.pdf}{
The complete integer census contains $15{,}198{,}743$ primitive triples. Exactly $121$ have $c>R$. Those $121$ points appear here. Classification uses exact integer comparison, while the displayed logarithms use floating point.\par
The highlighted point is $1+4374=4375$, with $R=210$ and quality approximately $1.567887$. It has the highest displayed quality in this census.\par
The line $c=R$ is a reference line. It is not the abc bound with its unknown constant. Points above it are expected exceptions to the exponent one comparison. No finite search determines a constant that works for every triple.
}{State the full search domain before interpreting a scatter plot.}

\GuideSpread{The exponent slack cannot be removed}{An unconditional family rules out $c\le KR$ uniformly}{05_epsilon_zero.pdf}{
This alternate family takes $(a,b,c)=(1,3^{2^n}-1,3^{2^n})$ for $n\ge1$. These triples are primitive. The lifting the exponent identity gives $v_2(b)=n+2$.\par
Because $\rad(b)\le b/2^{n+1}$ and $\rad(c)=3$, we get $R\le3b/2^{n+1}$. Consequently $c/R>2^{n+1}/3$, which tends to infinity. The chart displays the exact rational lower bound for the first six values.\par
A fixed constant cannot absorb this growth at exponent one. Positive $\varepsilon$ in abc is a mathematical necessity, rather than a technical decoration.
}{The $\varepsilon=0$ statement is false, independently of the unresolved conjecture.}

\GuideSpread{Uniformity is the main burden}{Choose $\varepsilon$, then its constant, then every triple}{06_quantifiers.pdf}{
The conjecture says that for every $\varepsilon>0$ there is $K_\varepsilon>0$ such that every primitive positive triple satisfies $c\le K_\varepsilon R^{1+\varepsilon}$. The constant may depend on $\varepsilon$. It may not depend on the triple or its prime set.\par
Fixing a finite support set already gives finiteness of the possible triples through known unconditional bounds in terms of the radical. This conclusion allows the support itself to control the constant.\par
The desired radical scale bound must hold while support and exponents vary together. Neither a fixed support computation nor a finite radical range supplies this quantifier order.
}{Check which variables every constant is allowed to depend on.}

\GuideSpread{The polynomial analogue has a working derivative}{Mason--Stothers in a sharp small example}{07_mason.pdf}{
Over a field of characteristic zero, nonzero pairwise coprime polynomials with $A+B=C$ and at least one nonconstant satisfy the polynomial abc bound. The radical counts distinct roots.\par
For $A=X^3$, $B=1-X^3$, and $C=1$, the radical has degree four and the maximum degree is three. The theorem gives $3\le4-1$, with equality.\par
The derivative Wronskian is $A'B-AB'=3X^2$. It sees the repeated root of $A$ and remains nonzero. The derivative and nonconstant hypotheses are load bearing. In positive characteristic the corresponding derivative condition must be checked explicitly.
}{The polynomial theorem works through a nonzero derivative, under its exact hypotheses.}

\GuideSpread{The integer setting needs another comparison}{A derivation on $\Z$ is forced to vanish}{08_integer_gap.pdf}{
An additive derivation obeying the Leibniz rule has $D(1)=0$. Therefore $D(n)=nD(1)=0$ for every integer $n$. Ordinary integer differentiation cannot reproduce the nonzero polynomial Wronskian.\par
The project instead constructs prime packet differentials with explicit weighted coordinates. Their local divisibility statements are valuable, but a small global transverse packet is not automatic.\par
The gap is a missing arithmetic comparison, rather than an absence of local formulas. This observation does not rule out other global methods, or turn a formal analogy into a proof of integer abc.
}{Preserve the mechanism of a theorem before transferring its conclusion to a new domain.}

\GuideSpread{The balanced defect corrects the target}{The factor $ab$ is essential}{09_balanced_defect.pdf}{
Write $Q=abc/R$. The exact identity is $Q/(ab)=c/R=M_c/(R_aR_b)$. Thus abc controls the balanced ratio after the input factor is removed.\par
For $1+8=9$, the cancellation is $12/8=3/2$. A theorem claiming $Q=R^{o(1)}$ uniformly would be a different statement. The recovered September block explicitly retracts that target.\par
The family $(1,2^k,2^k+1)$ has $Q\ge2^{k-1}$ and $R\le3\cdot2^k$. It rules out every proposed uniform bound $Q\le KR^\eta(\log(3c))^A$ with fixed $\eta<1$. The local jet identities survive this correction.
}{Normalize the arithmetic target before spending a local contact gain.}

\GuideSpread{The output sector gives an equivalent target}{Separate the input bandwidth from output persistence}{10_prime_sectors.pdf}{
Set $R_u=\rad(u)$ and $M_u=u/R_u$. For the high quality example, $M_a=1$, $M_b=3^9$, and $M_c=23^4$. The supports remain pairwise disjoint.\par
The bound $M_c\le K_\varepsilon R_aR_bR^\varepsilon$ is exactly the ordinary abc bound after multiplying by $R_c$. It keeps the requested output depth visible.\par
The inherited coherence route factors this output powerful part as $M_c=\epsilon_2KWC$. Here $K$ is a Kummer factor and $W$ is a logarithmic defect factor. That $W$ has a different meaning from the integer witness used at the end of this guide.
}{An exact factorization isolates the remaining arithmetic estimate without proving it.}

\GuideSpread{Finite cyclic controls distinguish zero from presence}{A complete toy model for the COS counting interface}{11_cyclic_controls.pdf}{
This figure enumerates all nine assignments $u,v\in\{-1,0,1\}$ in the cyclic group $\Z/13\Z$, with target residue $h=u+2v$. It is a toy model, rather than a recovered arithmetic channel.\par
There are exactly two assignments at $h=1$ and none at $h=6$. Cyclic Fourier orthogonality can recover these counts. Any upper certificate for a count must dominate the actual count, including the positive control.\par
The project's signed filtering route seeks a low complexity certificate for actual prime channels. A finite zero certificate does not supply a uniform arithmetic rigidity theorem, while a positive count cannot be eliminated by spectral rearrangement.
}{A useful filter must pass both positive and zero controls in its declared domain.}

\GuideSpread{A short relation does not control arbitrary index}{An exact lattice family exposes the missing arithmetic input}{12_arithmetic_gate.pdf}{
The lattice $\Lambda_N=\{(x,y)\in\Z^2:x-y\equiv0\pmod N\}$ has index $N$. The primitive vector $(1,1)$ belongs to it for every $N$. Its length stays fixed while the index grows without limit.\par
The coherence construction likewise has a normalized relation vector in a congruence kernel. A generic short vector theorem cannot bound its index from that vector alone.\par
What may distinguish the project lattice is the arithmetic origin of the rows: logarithms of actual input primes at actual output places, constrained by the additive equation. The missing theorem must use that structure quantitatively.
}{The arithmetic source of a row matters more than the mere existence of a short kernel vector.}

\GuideSpread{Root extraction has a geometric bill}{The full Fermat curve restores a factor of order $M$}{13_fermat_cost.pdf}{
Choose $u^M=a/c$ and $v^M=b/c$. The Fermat point $[u:v:1]$ has projective height $(\log c)/M$. This coefficient decreases as $M$ grows.\par
For the smooth plane Fermat curve, the canonical bundle is $\mathcal O(M-3)$. Its height coefficient is $(M-3)/M$. The bounded normalization term is for each fixed $M$, with no uniform bound asserted. The induced tensor power metric gives an exact coefficient equality.\par
The chart plots these exact rational coefficients for $4\le M\le24$. It concerns the full Fermat model and its line bundle. It does not exclude a suitably constructed sparse quotient with another metric. Such a quotient would need its own degree and norm accounting.
}{A coordinate height saving must survive the global line bundle and metric.}

\GuideSpread{The jet filter creates genuine Taylor cancellation}{Exact coefficients for the first five filters}{14_jet_series.pdf}{
For $m\ge0$, put $P_m(X)=X\prod_{j=1}^m(jX-1)$. Its coefficients, divided by $(m+1)!$, define the mixed weight polylogarithm combination $\Phi_m(t)$.\par
The exact series is $(-1)^m\sum_{k\ge1}\binom{k}{m+1}t^k/k^{m+2}$. Its first nonzero term is at $k=m+1$, with coefficient $(-1)^m/(m+1)^{m+2}$. The table displays that contact and the numerator coefficients.\par
The numerator vector is primitive and its normalized coefficient $\ell^1$ norm is one. This norm is a coefficient identity. It does not bound the height of a period realization, and the number of participating weights grows with $m$.
}{Formal contact order and arithmetic realization cost are separate invariants.}

\GuideSpread{A determinant retains independent moment directions}{The first nontrivial rank is two}{15_rank_two.pdf}{
Define $\mathcal H_r(t)=\det(\Li_{i+j+1}(t))_{0\le i,j<r}$ for $r\ge1$. At rank two it is $\Li_1(t)\Li_3(t)-\Li_2(t)^2$.\par
The moment matrix is a sum of rank one matrices with nodes $1/k$ and weights $t^k/k$. Cauchy--Binet turns the determinant into a sum over distinct nodes.\par
For rank two, the contribution from $i<j$ is $t^{i+j}(j-i)^2/(i^3j^3)$. The unique first pair is $\{1,2\}$, giving $t^3/8$. The cancellation has a concrete Taylor meaning, independent of any later arithmetic interpretation.
}{The determinant's nonzero initial term can be proved before constructing a global witness.}

\GuideSpread{Contact order grows quadratically with rank}{The unique initial subset is $\{1,\ldots,r\}$}{16_contact_order.pdf}{
For each rank, the Cauchy--Binet expansion has powers $t^{k_1+\cdots+k_r}$ over increasing positive indices. The unique minimum sum is $1+\cdots+r$.\par
Thus $s_r=r(r+1)/2$ is the exact contact order. The leading coefficient is $\prod_{k=1}^r k^{-(2k-1)}$, which is nonzero. The first five ranks have orders $1,3,6,10,15$.\par
The displayed fractions were computed exactly and checked against the subset expansion. This quadratic growth is useful local information. Turning it into a radical scale theorem still requires every place and every global cost to be compared within one arithmetic object.
}{A larger local contact multiplier is useful only if the comparison cost grows more slowly.}

\GuideSpread{Positive moments give real nonvanishing}{A finite rational example at $t=1/2$}{17_positive_moments.pdf}{
The bars are the fifteen pair contributions from indices $1\le i<j\le6$ to the rank two truncated determinant at $t=1/2$. Each contribution is an exact positive rational number.\par
For the full series, the positive Cauchy--Binet expansion proves $\mathcal H_r(t)>0$ when $0<t<1$. Hadamard and a convergent zeta product give $\mathcal H_r(t)\le e^{1/4}[-\log(1-t)]$.\par
For $a/c$ and $b/c$, this is a useful complex bound. It does not prove nonvanishing at the third chart $c/a>1$, or identify these complex values with the places of one algebraic number.
}{Use positivity only in its real interval of validity.}

\GuideSpread{Good primes give an exact universal valuation}{The strict threshold is $p>3r-1$}{18_good_primes.pdf}{
If $v_p(t)=e\ge1$ and $p>3r-1$, the undepleted series has exact valuation $v_p(\mathcal H_r(t))=es_r$. The first subset coefficient is a unit and every other subset has larger valuation.\par
The cards show three valid rank and prime pairs. The threshold is strict. At the boundary $r=2$, $p=5$, the independently verified argument $t=10$ has raw valuation four, exceeding the contact value three.\par
At an output prime dividing $c$, the argument $c/a$ is small $p$ adically because $p\nmid a$. The argument $a/c$ is not small there. Chart and place must be chosen together.
}{Exact local equality depends on the prime threshold and the valuation of the selected argument.}

\GuideSpread{Depletion changes the first surviving index}{A local function depends on the chosen prime}{19_depletion.pdf}{
The depleted jet series deletes terms with $p\mid k$. When $p\nmid m+1$, its first index remains $m+1$. When $p\mid m+1$, the first surviving index becomes $m+2$.\par
In the example $m=3$ and $p=2$, index four is removed and index five survives. Its coefficient is $-1/625$, a $2$ adic unit. For $v_2(t)=e\ge1$, the exact valuation is $5e$.\par
This improves local integrality, but it introduces dependence on $p$. The various depleted functions are not automatically the local values of one nonzero algebraic witness. A product formula needs that additional identification.
}{A prime dependent filter supplies a local theorem and creates a separate realization question.}

\GuideSpread{Residue collisions add exact valuation depth}{Unit index sums and Vandermonde energy}{20_residue_energy.pdf}{
Let $u_1<\cdots<u_r$ be the first positive indices not divisible by $p$. Put $S_{p,r}=\sum u_i$ and let $E_{p,r}$ be the valuation of their Vandermonde product.\par
The first unit indices simultaneously minimize the index sum and the residue collision energy at every power of $p$. Their term has unique minimum valuation. Therefore $v_p(\mathcal H_{r,p}(t))=eS_{p,r}+2E_{p,r}$ for $e=v_p(t)\ge1$.\par
At rank four and $e=1$, the valuations are $32$, $16$, and $10$ for $p=2,3,5$. These exact integers were derived from the indices and matched the fresh replay.
}{Small prime residue collisions contribute a measurable local correction.}

\GuideSpread{A finite budget controls the raw small primes}{Rank two has the safe correction $B_2=3456$}{21_small_prime_budget.pdf}{
The raw universal series has denominators at small primes. The recovered proof supplies a fixed integer $B_r$, depending only on rank, with $v_p(B_r\mathcal H_r(t))\ge es_r$ whenever $v_p(t)=e\ge1$.\par
For rank two, one valid correction is $2^7 3^3=3456$. The figure shows these budget exponents. It does not claim that the correction is minimal.\par
The small prime bound permits $\mathcal H_r(t)=0$, with $v_p(0)=+\infty$. A zero correction exponent also does not imply equality of valuation. This is a theorem about local convergent analytic series, rather than the denominator of a constructed motivic or algebraic class.
}{A finite denominator budget is useful evidence, but it is not an all place arithmetic identity.}

\GuideSpread{The simplest scalar candidate fails a contact gate}{Exact truncation loses output prime divisibility}{22_truncation_loss.pdf}{
Truncate the rank two determinant at index three. Its exact polynomial is $\mathcal H_{2,3}(t)=t^3(27+32t+t^2)/216$. Clearing denominators at $t=a/c$ gives the integer $W=a^3(27c^2+32ac+a^2)$.\par
If $p>3$ divides $c$ and $p\nmid a$, then $W\equiv a^5\not\equiv0\pmod p$. This scalar carries no positive output contact at that prime.\par
The general truncation theorem has the same obstruction at output primes $p>K$. The candidate is genuinely algebraic, but it fails the required local divisibility. This does not rule out every determinant construction.
}{An arithmetic witness must preserve the required contact after clearing all denominators.}

\GuideSpread{The status record has several different meanings}{Primary source checks through 8 October 2026}{23_evidence_status.pdf}{
The IUT papers were published in the RIMS special issue in 2021. Publication authenticates a publication record. It is separate from resolving the disputed implication in the proof claim.\par
The July 2026 Project LANA report suspends judgment on whether IUT proves abc. It does not reconstruct the required comparison in the transition from IUT III Theorem 3.11 to Corollary 3.12. The monograph cites that primary report and the competing author responses.\par
This guide does not audit the full IUT argument or survey every mathematician. Our exact local results and bounded computations also leave their stated global gates unresolved. A proof claim, a publication, and a finite replay are distinct evidence types.
}{Read each source for the claim it actually supports, with its date and scope retained.}

\GuideSpread{The global bridge states the remaining burden}{A conditional certificate with the correct input scale}{24_global_bridge.pdf}{
For selected ranks, suppose every primitive triple admits a nonzero integer witness $W$ with $Q^{s_r}\mid W$ and $|W|\le C_r(ab)^{s_r}R^{E_r}(\log(3c))^{A_r}$. Constants may depend on rank, but not on the triple.\par
Divisibility and nonvanishing give $Q^{s_r}\le|W|$. Dividing by $(ab)^{s_r}$ produces the balanced estimate in the figure. If $E_r/s_r\to0$, logarithmic powers can be absorbed into an arbitrary additional radical exponent, yielding abc.\par
The existence, nonvanishing, full local divisibility, and global height bound of this witness remain unproved for the determinant candidate. The dashed arrows mark those hypotheses.
}{The local contact theorem must meet a correctly normalized, nonzero global comparison.}
